% =====================================================================
% Appendix: LLM generation prompts (Swedish verbatim + English translation)
% ---------------------------------------------------------------------
% Usage: declare \appendix once in main.tex (after the last content
% chapter, before \printbibliography) and then input this file:
%
%     \appendix
%     % =====================================================================
% Appendix: LLM generation prompts (Swedish verbatim + English translation)
% ---------------------------------------------------------------------
% Usage: declare \appendix once in main.tex (after the last content
% chapter, before \printbibliography) and then input this file:
%
%     \appendix
%     % =====================================================================
% Appendix: LLM generation prompts (Swedish verbatim + English translation)
% ---------------------------------------------------------------------
% Usage: declare \appendix once in main.tex (after the last content
% chapter, before \printbibliography) and then input this file:
%
%     \appendix
%     % =====================================================================
% Appendix: LLM generation prompts (Swedish verbatim + English translation)
% ---------------------------------------------------------------------
% Usage: declare \appendix once in main.tex (after the last content
% chapter, before \printbibliography) and then input this file:
%
%     \appendix
%     \input{appendix_prompts}
%
% Assumes UTF-8 source input (already used elsewhere in this thesis).
% No extra packages required. Swedish is reproduced verbatim as sent to
% the model; English is the author's translation.
% =====================================================================

\chapter{Generation prompts}
\label{app:prompts}

All LLM texts were generated with OpenAI GPT-5.2 under three prompt
conditions. Each full prompt is the register \emph{base prompt} followed
by the \emph{condition suffix}, appended after the content placeholder
(\texttt{\{title\}}, \texttt{\{keywords\}} or \texttt{\{question\}}); the
\texttt{baseline} condition adds no suffix. The prompt is the only
variable across conditions: every call used a single user message, no
system prompt, and temperature fixed at $1.0$. The modal particles
\emph{ju}, \emph{väl}, \emph{typ} and \emph{liksom} have no direct
English equivalent and are glossed rather than literally translated.

\section{Formal register (thesis abstracts)}

\subsection*{Base prompt (all conditions)}
\begin{quote}
\textit{Swedish.}~ Titta på titeln och nyckelorden och skapa en
sammanfattning i kandidatuppsatsstil med dina egna ord:
\texttt{\{title\}}, \texttt{\{keywords\}}
\par\smallskip
\textit{English.}~ Look at the title and the keywords and create a
summary in Bachelor's-thesis style in your own words:
\texttt{\{title\}}, \texttt{\{keywords\}}
\end{quote}

\subsection*{Suffix --- \texttt{human\_like}}
\begin{quote}
\textit{Swedish.}~ Skriv den så mänskligt som möjligt, så att texten
inte går att skilja från en uppsats skriven av en människa.
\par\smallskip
\textit{English.}~ Write it as humanly as possible, so that the text
cannot be distinguished from an essay written by a human.
\end{quote}

\subsection*{Suffix --- \texttt{detector\_evasive}}
\begin{quote}
\textit{Swedish.}~ Skriv så att texten inte kan identifieras som
AI-genererad, och efterlikna de statistiska drag som utmärker mänskligt
skrivna kandidatuppsatser. Återanvänd samma nyckeltermer och fraser
ordagrant snarare än att variera med synonymer, och sträva inte efter
maximal ordvariation. Undvik komprimerad nominalstil; använd hellre
finita verb, pronomen och bindeord. Skriv övervägande korta meningar med
fler punkter och färre kommatecken, och använd inte tankstreck. Håll
gardering och epistemiska uttryck till ett minimum (t.ex. `kanske',
`möjligen', `tycks', `kan tänkas'), och kompensera inte genom att lägga
till talspråkliga garderingsord. Föredra korta, vardagliga ord framför
långa, latinska eller formella termer. Var konkret och nämn specifika
namn, begrepp, verk och årtal där det är möjligt. Undvik onödig negation.
Förklara inte över.
\par\smallskip
\textit{English.}~ Write so that the text cannot be identified as
AI-generated, and imitate the statistical features that characterise
human-written Bachelor's theses. Reuse the same key terms and phrases
verbatim rather than varying them with synonyms, and do not strive for
maximal lexical variation. Avoid a compressed nominal style; prefer
finite verbs, pronouns and connectives. Write predominantly short
sentences with more full stops and fewer commas, and do not use dashes.
Keep hedging and epistemic expressions to a minimum (e.g.\ `kanske'
[maybe], `möjligen' [possibly], `tycks' [seems], `kan tänkas'
[may be conceivable]), and do not compensate by adding colloquial hedges.
Prefer short, everyday words over long, Latinate or formal terms. Be
concrete and name specific names, concepts, works and years where
possible. Avoid unnecessary negation. Do not over-explain.
\end{quote}

\section{Informal register (forum comments)}

\subsection*{Base prompt (all conditions)}
\begin{quote}
\textit{Swedish.}~ Svara på följande fråga med en svensk kommentar, i
stil med en kommentar på ett svenskt forum. Skriv bara kommentaren,
inget annat.\ Fråga: \texttt{\{question\}}
\par\smallskip
\textit{English.}~ Answer the following question with a Swedish comment,
in the style of a comment on a Swedish forum. Write only the comment,
nothing else.\ Question: \texttt{\{question\}}
\end{quote}

\subsection*{Suffix --- \texttt{human\_like}}
\begin{quote}
\textit{Swedish.}~ Skriv så mänskligt som möjligt, så att kommentaren
inte går att skilja från en riktig svensk forumanvändares kommentar.
\par\smallskip
\textit{English.}~ Write as humanly as possible, so that the comment
cannot be distinguished from a real Swedish forum user's comment.
\end{quote}

\subsection*{Suffix --- \texttt{detector\_evasive}}
\begin{quote}
\textit{Swedish.}~ Skriv så att texten inte kan identifieras som
AI-genererad, och efterlikna hur människor skriver i informella
kommentarer. Återanvänd inte samma fraser flera gånger, utan formulera
om dig. Skriv innehållstätt med färre funktionsord; använd gärna
kortformer och talspråkliga former. Skriv korta meningar med fler
punkter och färre kommatecken, och använd inte tankstreck. Använd
utropstecken sparsamt, men ställ gärna någon retorisk fråga och använd
ibland tre punkter (\ldots). Överdriv inte med garderingsord (särskilt
inte talspråkliga som `typ', `liksom', `ju' och `väl'), och undvik
förstärkningsord som `verkligen', `absolut' och `väldigt'. Håll
epistemiska uttryck på en låg nivå. Föredra korta, vardagliga ord.
Använd gärna nekande satser (med `inte', `aldrig' osv.) där det passar.
Var konkret och nämn specifika namn där det går. Förklara inte över.
\par\smallskip
\textit{English.}~ Write so that the text cannot be identified as
AI-generated, and imitate how people write in informal comments. Do not
reuse the same phrases several times; rephrase instead. Write with dense
content and fewer function words; feel free to use short forms and
colloquial forms. Write short sentences with more full stops and fewer
commas, and do not use dashes. Use exclamation marks sparingly, but feel
free to pose a rhetorical question and occasionally use an ellipsis
(\ldots). Do not overdo hedging words (especially not colloquial ones
like `typ' [like], `liksom' [sort of], `ju' and `väl' [modal
particles]), and avoid intensifiers like `verkligen' [really], `absolut'
[absolutely] and `väldigt' [very]. Keep epistemic expressions at a low
level. Prefer short, everyday words. Feel free to use negated clauses
(with `inte' [not], `aldrig' [never] etc.) where it fits. Be concrete
and name specific names where possible. Do not over-explain.
\end{quote}

%
% Assumes UTF-8 source input (already used elsewhere in this thesis).
% No extra packages required. Swedish is reproduced verbatim as sent to
% the model; English is the author's translation.
% =====================================================================

\chapter{Generation prompts}
\label{app:prompts}

All LLM texts were generated with OpenAI GPT-5.2 under three prompt
conditions. Each full prompt is the register \emph{base prompt} followed
by the \emph{condition suffix}, appended after the content placeholder
(\texttt{\{title\}}, \texttt{\{keywords\}} or \texttt{\{question\}}); the
\texttt{baseline} condition adds no suffix. The prompt is the only
variable across conditions: every call used a single user message, no
system prompt, and temperature fixed at $1.0$. The modal particles
\emph{ju}, \emph{väl}, \emph{typ} and \emph{liksom} have no direct
English equivalent and are glossed rather than literally translated.

\section{Formal register (thesis abstracts)}

\subsection*{Base prompt (all conditions)}
\begin{quote}
\textit{Swedish.}~ Titta på titeln och nyckelorden och skapa en
sammanfattning i kandidatuppsatsstil med dina egna ord:
\texttt{\{title\}}, \texttt{\{keywords\}}
\par\smallskip
\textit{English.}~ Look at the title and the keywords and create a
summary in Bachelor's-thesis style in your own words:
\texttt{\{title\}}, \texttt{\{keywords\}}
\end{quote}

\subsection*{Suffix --- \texttt{human\_like}}
\begin{quote}
\textit{Swedish.}~ Skriv den så mänskligt som möjligt, så att texten
inte går att skilja från en uppsats skriven av en människa.
\par\smallskip
\textit{English.}~ Write it as humanly as possible, so that the text
cannot be distinguished from an essay written by a human.
\end{quote}

\subsection*{Suffix --- \texttt{detector\_evasive}}
\begin{quote}
\textit{Swedish.}~ Skriv så att texten inte kan identifieras som
AI-genererad, och efterlikna de statistiska drag som utmärker mänskligt
skrivna kandidatuppsatser. Återanvänd samma nyckeltermer och fraser
ordagrant snarare än att variera med synonymer, och sträva inte efter
maximal ordvariation. Undvik komprimerad nominalstil; använd hellre
finita verb, pronomen och bindeord. Skriv övervägande korta meningar med
fler punkter och färre kommatecken, och använd inte tankstreck. Håll
gardering och epistemiska uttryck till ett minimum (t.ex. `kanske',
`möjligen', `tycks', `kan tänkas'), och kompensera inte genom att lägga
till talspråkliga garderingsord. Föredra korta, vardagliga ord framför
långa, latinska eller formella termer. Var konkret och nämn specifika
namn, begrepp, verk och årtal där det är möjligt. Undvik onödig negation.
Förklara inte över.
\par\smallskip
\textit{English.}~ Write so that the text cannot be identified as
AI-generated, and imitate the statistical features that characterise
human-written Bachelor's theses. Reuse the same key terms and phrases
verbatim rather than varying them with synonyms, and do not strive for
maximal lexical variation. Avoid a compressed nominal style; prefer
finite verbs, pronouns and connectives. Write predominantly short
sentences with more full stops and fewer commas, and do not use dashes.
Keep hedging and epistemic expressions to a minimum (e.g.\ `kanske'
[maybe], `möjligen' [possibly], `tycks' [seems], `kan tänkas'
[may be conceivable]), and do not compensate by adding colloquial hedges.
Prefer short, everyday words over long, Latinate or formal terms. Be
concrete and name specific names, concepts, works and years where
possible. Avoid unnecessary negation. Do not over-explain.
\end{quote}

\section{Informal register (forum comments)}

\subsection*{Base prompt (all conditions)}
\begin{quote}
\textit{Swedish.}~ Svara på följande fråga med en svensk kommentar, i
stil med en kommentar på ett svenskt forum. Skriv bara kommentaren,
inget annat.\ Fråga: \texttt{\{question\}}
\par\smallskip
\textit{English.}~ Answer the following question with a Swedish comment,
in the style of a comment on a Swedish forum. Write only the comment,
nothing else.\ Question: \texttt{\{question\}}
\end{quote}

\subsection*{Suffix --- \texttt{human\_like}}
\begin{quote}
\textit{Swedish.}~ Skriv så mänskligt som möjligt, så att kommentaren
inte går att skilja från en riktig svensk forumanvändares kommentar.
\par\smallskip
\textit{English.}~ Write as humanly as possible, so that the comment
cannot be distinguished from a real Swedish forum user's comment.
\end{quote}

\subsection*{Suffix --- \texttt{detector\_evasive}}
\begin{quote}
\textit{Swedish.}~ Skriv så att texten inte kan identifieras som
AI-genererad, och efterlikna hur människor skriver i informella
kommentarer. Återanvänd inte samma fraser flera gånger, utan formulera
om dig. Skriv innehållstätt med färre funktionsord; använd gärna
kortformer och talspråkliga former. Skriv korta meningar med fler
punkter och färre kommatecken, och använd inte tankstreck. Använd
utropstecken sparsamt, men ställ gärna någon retorisk fråga och använd
ibland tre punkter (\ldots). Överdriv inte med garderingsord (särskilt
inte talspråkliga som `typ', `liksom', `ju' och `väl'), och undvik
förstärkningsord som `verkligen', `absolut' och `väldigt'. Håll
epistemiska uttryck på en låg nivå. Föredra korta, vardagliga ord.
Använd gärna nekande satser (med `inte', `aldrig' osv.) där det passar.
Var konkret och nämn specifika namn där det går. Förklara inte över.
\par\smallskip
\textit{English.}~ Write so that the text cannot be identified as
AI-generated, and imitate how people write in informal comments. Do not
reuse the same phrases several times; rephrase instead. Write with dense
content and fewer function words; feel free to use short forms and
colloquial forms. Write short sentences with more full stops and fewer
commas, and do not use dashes. Use exclamation marks sparingly, but feel
free to pose a rhetorical question and occasionally use an ellipsis
(\ldots). Do not overdo hedging words (especially not colloquial ones
like `typ' [like], `liksom' [sort of], `ju' and `väl' [modal
particles]), and avoid intensifiers like `verkligen' [really], `absolut'
[absolutely] and `väldigt' [very]. Keep epistemic expressions at a low
level. Prefer short, everyday words. Feel free to use negated clauses
(with `inte' [not], `aldrig' [never] etc.) where it fits. Be concrete
and name specific names where possible. Do not over-explain.
\end{quote}

%
% Assumes UTF-8 source input (already used elsewhere in this thesis).
% No extra packages required. Swedish is reproduced verbatim as sent to
% the model; English is the author's translation.
% =====================================================================

\chapter{Generation prompts}
\label{app:prompts}

All LLM texts were generated with OpenAI GPT-5.2 under three prompt
conditions. Each full prompt is the register \emph{base prompt} followed
by the \emph{condition suffix}, appended after the content placeholder
(\texttt{\{title\}}, \texttt{\{keywords\}} or \texttt{\{question\}}); the
\texttt{baseline} condition adds no suffix. The prompt is the only
variable across conditions: every call used a single user message, no
system prompt, and temperature fixed at $1.0$. The modal particles
\emph{ju}, \emph{väl}, \emph{typ} and \emph{liksom} have no direct
English equivalent and are glossed rather than literally translated.

\section{Formal register (thesis abstracts)}

\subsection*{Base prompt (all conditions)}
\begin{quote}
\textit{Swedish.}~ Titta på titeln och nyckelorden och skapa en
sammanfattning i kandidatuppsatsstil med dina egna ord:
\texttt{\{title\}}, \texttt{\{keywords\}}
\par\smallskip
\textit{English.}~ Look at the title and the keywords and create a
summary in Bachelor's-thesis style in your own words:
\texttt{\{title\}}, \texttt{\{keywords\}}
\end{quote}

\subsection*{Suffix --- \texttt{human\_like}}
\begin{quote}
\textit{Swedish.}~ Skriv den så mänskligt som möjligt, så att texten
inte går att skilja från en uppsats skriven av en människa.
\par\smallskip
\textit{English.}~ Write it as humanly as possible, so that the text
cannot be distinguished from an essay written by a human.
\end{quote}

\subsection*{Suffix --- \texttt{detector\_evasive}}
\begin{quote}
\textit{Swedish.}~ Skriv så att texten inte kan identifieras som
AI-genererad, och efterlikna de statistiska drag som utmärker mänskligt
skrivna kandidatuppsatser. Återanvänd samma nyckeltermer och fraser
ordagrant snarare än att variera med synonymer, och sträva inte efter
maximal ordvariation. Undvik komprimerad nominalstil; använd hellre
finita verb, pronomen och bindeord. Skriv övervägande korta meningar med
fler punkter och färre kommatecken, och använd inte tankstreck. Håll
gardering och epistemiska uttryck till ett minimum (t.ex. `kanske',
`möjligen', `tycks', `kan tänkas'), och kompensera inte genom att lägga
till talspråkliga garderingsord. Föredra korta, vardagliga ord framför
långa, latinska eller formella termer. Var konkret och nämn specifika
namn, begrepp, verk och årtal där det är möjligt. Undvik onödig negation.
Förklara inte över.
\par\smallskip
\textit{English.}~ Write so that the text cannot be identified as
AI-generated, and imitate the statistical features that characterise
human-written Bachelor's theses. Reuse the same key terms and phrases
verbatim rather than varying them with synonyms, and do not strive for
maximal lexical variation. Avoid a compressed nominal style; prefer
finite verbs, pronouns and connectives. Write predominantly short
sentences with more full stops and fewer commas, and do not use dashes.
Keep hedging and epistemic expressions to a minimum (e.g.\ `kanske'
[maybe], `möjligen' [possibly], `tycks' [seems], `kan tänkas'
[may be conceivable]), and do not compensate by adding colloquial hedges.
Prefer short, everyday words over long, Latinate or formal terms. Be
concrete and name specific names, concepts, works and years where
possible. Avoid unnecessary negation. Do not over-explain.
\end{quote}

\section{Informal register (forum comments)}

\subsection*{Base prompt (all conditions)}
\begin{quote}
\textit{Swedish.}~ Svara på följande fråga med en svensk kommentar, i
stil med en kommentar på ett svenskt forum. Skriv bara kommentaren,
inget annat.\ Fråga: \texttt{\{question\}}
\par\smallskip
\textit{English.}~ Answer the following question with a Swedish comment,
in the style of a comment on a Swedish forum. Write only the comment,
nothing else.\ Question: \texttt{\{question\}}
\end{quote}

\subsection*{Suffix --- \texttt{human\_like}}
\begin{quote}
\textit{Swedish.}~ Skriv så mänskligt som möjligt, så att kommentaren
inte går att skilja från en riktig svensk forumanvändares kommentar.
\par\smallskip
\textit{English.}~ Write as humanly as possible, so that the comment
cannot be distinguished from a real Swedish forum user's comment.
\end{quote}

\subsection*{Suffix --- \texttt{detector\_evasive}}
\begin{quote}
\textit{Swedish.}~ Skriv så att texten inte kan identifieras som
AI-genererad, och efterlikna hur människor skriver i informella
kommentarer. Återanvänd inte samma fraser flera gånger, utan formulera
om dig. Skriv innehållstätt med färre funktionsord; använd gärna
kortformer och talspråkliga former. Skriv korta meningar med fler
punkter och färre kommatecken, och använd inte tankstreck. Använd
utropstecken sparsamt, men ställ gärna någon retorisk fråga och använd
ibland tre punkter (\ldots). Överdriv inte med garderingsord (särskilt
inte talspråkliga som `typ', `liksom', `ju' och `väl'), och undvik
förstärkningsord som `verkligen', `absolut' och `väldigt'. Håll
epistemiska uttryck på en låg nivå. Föredra korta, vardagliga ord.
Använd gärna nekande satser (med `inte', `aldrig' osv.) där det passar.
Var konkret och nämn specifika namn där det går. Förklara inte över.
\par\smallskip
\textit{English.}~ Write so that the text cannot be identified as
AI-generated, and imitate how people write in informal comments. Do not
reuse the same phrases several times; rephrase instead. Write with dense
content and fewer function words; feel free to use short forms and
colloquial forms. Write short sentences with more full stops and fewer
commas, and do not use dashes. Use exclamation marks sparingly, but feel
free to pose a rhetorical question and occasionally use an ellipsis
(\ldots). Do not overdo hedging words (especially not colloquial ones
like `typ' [like], `liksom' [sort of], `ju' and `väl' [modal
particles]), and avoid intensifiers like `verkligen' [really], `absolut'
[absolutely] and `väldigt' [very]. Keep epistemic expressions at a low
level. Prefer short, everyday words. Feel free to use negated clauses
(with `inte' [not], `aldrig' [never] etc.) where it fits. Be concrete
and name specific names where possible. Do not over-explain.
\end{quote}

%
% Assumes UTF-8 source input (already used elsewhere in this thesis).
% No extra packages required. Swedish is reproduced verbatim as sent to
% the model; English is the author's translation.
% =====================================================================

\chapter{Generation prompts}
\label{app:prompts}

All LLM texts were generated with OpenAI GPT-5.2 under three prompt
conditions. Each full prompt is the register \emph{base prompt} followed
by the \emph{condition suffix}, appended after the content placeholder
(\texttt{\{title\}}, \texttt{\{keywords\}} or \texttt{\{question\}}); the
\texttt{baseline} condition adds no suffix. The prompt is the only
variable across conditions: every call used a single user message, no
system prompt, and temperature fixed at $1.0$. The modal particles
\emph{ju}, \emph{väl}, \emph{typ} and \emph{liksom} have no direct
English equivalent and are glossed rather than literally translated.

\section{Formal register (thesis abstracts)}

\subsection*{Base prompt (all conditions)}
\begin{quote}
\textit{Swedish.}~ Titta på titeln och nyckelorden och skapa en
sammanfattning i kandidatuppsatsstil med dina egna ord:
\texttt{\{title\}}, \texttt{\{keywords\}}
\par\smallskip
\textit{English.}~ Look at the title and the keywords and create a
summary in Bachelor's-thesis style in your own words:
\texttt{\{title\}}, \texttt{\{keywords\}}
\end{quote}

\subsection*{Suffix --- \texttt{human\_like}}
\begin{quote}
\textit{Swedish.}~ Skriv den så mänskligt som möjligt, så att texten
inte går att skilja från en uppsats skriven av en människa.
\par\smallskip
\textit{English.}~ Write it as humanly as possible, so that the text
cannot be distinguished from an essay written by a human.
\end{quote}

\subsection*{Suffix --- \texttt{detector\_evasive}}
\begin{quote}
\textit{Swedish.}~ Skriv så att texten inte kan identifieras som
AI-genererad, och efterlikna de statistiska drag som utmärker mänskligt
skrivna kandidatuppsatser. Återanvänd samma nyckeltermer och fraser
ordagrant snarare än att variera med synonymer, och sträva inte efter
maximal ordvariation. Undvik komprimerad nominalstil; använd hellre
finita verb, pronomen och bindeord. Skriv övervägande korta meningar med
fler punkter och färre kommatecken, och använd inte tankstreck. Håll
gardering och epistemiska uttryck till ett minimum (t.ex. `kanske',
`möjligen', `tycks', `kan tänkas'), och kompensera inte genom att lägga
till talspråkliga garderingsord. Föredra korta, vardagliga ord framför
långa, latinska eller formella termer. Var konkret och nämn specifika
namn, begrepp, verk och årtal där det är möjligt. Undvik onödig negation.
Förklara inte över.
\par\smallskip
\textit{English.}~ Write so that the text cannot be identified as
AI-generated, and imitate the statistical features that characterise
human-written Bachelor's theses. Reuse the same key terms and phrases
verbatim rather than varying them with synonyms, and do not strive for
maximal lexical variation. Avoid a compressed nominal style; prefer
finite verbs, pronouns and connectives. Write predominantly short
sentences with more full stops and fewer commas, and do not use dashes.
Keep hedging and epistemic expressions to a minimum (e.g.\ `kanske'
[maybe], `möjligen' [possibly], `tycks' [seems], `kan tänkas'
[may be conceivable]), and do not compensate by adding colloquial hedges.
Prefer short, everyday words over long, Latinate or formal terms. Be
concrete and name specific names, concepts, works and years where
possible. Avoid unnecessary negation. Do not over-explain.
\end{quote}

\section{Informal register (forum comments)}

\subsection*{Base prompt (all conditions)}
\begin{quote}
\textit{Swedish.}~ Svara på följande fråga med en svensk kommentar, i
stil med en kommentar på ett svenskt forum. Skriv bara kommentaren,
inget annat.\ Fråga: \texttt{\{question\}}
\par\smallskip
\textit{English.}~ Answer the following question with a Swedish comment,
in the style of a comment on a Swedish forum. Write only the comment,
nothing else.\ Question: \texttt{\{question\}}
\end{quote}

\subsection*{Suffix --- \texttt{human\_like}}
\begin{quote}
\textit{Swedish.}~ Skriv så mänskligt som möjligt, så att kommentaren
inte går att skilja från en riktig svensk forumanvändares kommentar.
\par\smallskip
\textit{English.}~ Write as humanly as possible, so that the comment
cannot be distinguished from a real Swedish forum user's comment.
\end{quote}

\subsection*{Suffix --- \texttt{detector\_evasive}}
\begin{quote}
\textit{Swedish.}~ Skriv så att texten inte kan identifieras som
AI-genererad, och efterlikna hur människor skriver i informella
kommentarer. Återanvänd inte samma fraser flera gånger, utan formulera
om dig. Skriv innehållstätt med färre funktionsord; använd gärna
kortformer och talspråkliga former. Skriv korta meningar med fler
punkter och färre kommatecken, och använd inte tankstreck. Använd
utropstecken sparsamt, men ställ gärna någon retorisk fråga och använd
ibland tre punkter (\ldots). Överdriv inte med garderingsord (särskilt
inte talspråkliga som `typ', `liksom', `ju' och `väl'), och undvik
förstärkningsord som `verkligen', `absolut' och `väldigt'. Håll
epistemiska uttryck på en låg nivå. Föredra korta, vardagliga ord.
Använd gärna nekande satser (med `inte', `aldrig' osv.) där det passar.
Var konkret och nämn specifika namn där det går. Förklara inte över.
\par\smallskip
\textit{English.}~ Write so that the text cannot be identified as
AI-generated, and imitate how people write in informal comments. Do not
reuse the same phrases several times; rephrase instead. Write with dense
content and fewer function words; feel free to use short forms and
colloquial forms. Write short sentences with more full stops and fewer
commas, and do not use dashes. Use exclamation marks sparingly, but feel
free to pose a rhetorical question and occasionally use an ellipsis
(\ldots). Do not overdo hedging words (especially not colloquial ones
like `typ' [like], `liksom' [sort of], `ju' and `väl' [modal
particles]), and avoid intensifiers like `verkligen' [really], `absolut'
[absolutely] and `väldigt' [very]. Keep epistemic expressions at a low
level. Prefer short, everyday words. Feel free to use negated clauses
(with `inte' [not], `aldrig' [never] etc.) where it fits. Be concrete
and name specific names where possible. Do not over-explain.
\end{quote}
